\chapter{Making measurements projective}
\label{chap:projective}
\label{sec:making-measurements-projective}

For projective measurements, consistency and the state-dependent distance agree
up to the factor two in Proposition~\ref{prop:simeq-to-approx}.  The final
theorem therefore rounds the measurements produced by the induction to
projective measurements.  The relevant parameter is the state-dependent
idempotence defect $\sum_a\varphi(M_a-M_a^2)$, which is rounded with a linear
loss.  The lemma below is the finite-dimensional case of the orthogonalization
theorem of de~la~Salle \cite{deLaSalle2022Orthogonalization}.

\begin{definition}[Idempotence defect]\label{def:idempotence-defect}
	\lean{MIPStarRE.LDT.MakingMeasurementsProjective.SimplifiedOrthogonalization.idempotenceDefect}
	\leanok
	\uses{def:measurements}
	For a state $\varphi$ on $B(\calH)$ and a measurement $M=\{M_a\}$, set
	$\Delta=\sum_a\varphi(M_a-M_a^2)=1-\varphi(\sum_aM_a^2)$.
\end{definition}

\begin{lemma}[Global rank allocation]\label{lem:global-rank-allocation}
	\lean{MIPStarRE.LDT.MakingMeasurementsProjective.SimplifiedOrthogonalization.largest_weights_dominate_fractional, MIPStarRE.LDT.MakingMeasurementsProjective.SimplifiedOrthogonalization.exists_large_subset_ordered, MIPStarRE.LDT.MakingMeasurementsProjective.SimplifiedOrthogonalization.exists_selected_overlap_ge_one_sub_defect}
	\leanok
	\uses{def:idempotence-defect}
	Let $\calH$ have dimension $d$ and let $M_a=\sum_j\lambda_{a,j}\ket{v_{a,j}}\!\bra{v_{a,j}}$
	be spectral decompositions.  Put $w_{a,j}=\lambda_{a,j}\varphi(\ket{v_{a,j}}\!\bra{v_{a,j}})$.
	Selecting the $d$ largest weights and letting $Q_a$ project onto the
	selected eigenvectors of $M_a$ gives projections with
	\[
		[Q_a,M_a]=0,\qquad\sum_a\operatorname{rank}Q_a=d,\qquad
		\sum_a\varphi(Q_aM_a)\ge1-\Delta.
	\]
\end{lemma}
\begin{proof}
	\leanok
	Taking the trace of $\sum_aM_a=I$ gives $\sum_{a,j}\lambda_{a,j}=d$.  The
	sum of the $d$ largest weights dominates $\sum x_{a,j}w_{a,j}$ for every
	fractional selection $0\le x_{a,j}\le1$ of total mass $d$, in particular
	for $x_{a,j}=\lambda_{a,j}$, whose value is
	$\sum_a\varphi(M_a^2)=1-\Delta$.
\end{proof}

\begin{lemma}[Unitary polar factor]\label{lem:unitary-polar-factor}
	\lean{MIPStarRE.LDT.MakingMeasurementsProjective.SimplifiedOrthogonalization.exists_unitary_polar_factor}
	\leanok
	Every square complex matrix $X$ has a unitary $U$ with
	$X=U(X^\dagger X)^{1/2}$, also when $X$ is singular.
\end{lemma}
\begin{proof}
	\leanok
	Let $T=(X^\dagger X)^{1/2}$.  Since $\|Tv\|=\|Xv\|$ for every vector $v$,
	the assignment $Tv\mapsto Xv$ is a well-defined linear isometry on the
	range of $T$.  In finite dimension it extends to an isometry of the whole
	space, which is the required unitary.
\end{proof}

\begin{lemma}[State-dependent orthogonalization]\label{lem:state-dependent-orthogonalization}
	\lean{MIPStarRE.LDT.MakingMeasurementsProjective.SimplifiedOrthogonalization.exists_projective_measurement_linear_bound}
	\leanok
	\uses{def:idempotence-defect}
	Let $\calH$ be finite-dimensional, let $\varphi$ be a state on $B(\calH)$,
	and let $M=\{M_a\}$ be a measurement.  Then there is a projective
	measurement $P=\{P_a\}$ such that
	\[
		\sum_a\varphi\bigl((M_a-P_a)^2\bigr)\le9\Delta.
	\]
\end{lemma}
\begin{proof}
	\leanok
	\uses{lem:global-rank-allocation, lem:unitary-polar-factor}
	Take $Q_a$ from Lemma~\ref{lem:global-rank-allocation}, let
	$\calK=\bigoplus_aQ_a\calH$ with coordinate projections $E_a$, and define
	$X\colon\calH\to\calK$ by $Xv=(Q_aM_a^{1/2}v)_a$.  Then
	$X^\dagger X=\sum_aQ_aM_a=:S\le I$.  Write $\abs X=S^{1/2}$ and choose a
	unitary $U$ with $X=U\abs X$ by Lemma~\ref{lem:unitary-polar-factor}; this
	is possible because $\dim\calK=\sum_a\operatorname{rank}Q_a=\dim\calH$.  Set
	$P_a=U^\dagger E_aU$.  Then $\abs XP_a\abs X=Q_aM_a$ and
	\[
		M_a-P_a=(I-Q_a)M_a+(\abs X-I)P_a\abs X+P_a(\abs X-I).
	\]
	In the seminorm $\|Y\|_{\varphi,2}^2=\sum_a\varphi(Y_a^\dagger Y_a)$ each of
	the three families has squared norm at most $\Delta$: the first by
	$[Q_a,M_a]=0$ and $M_a^2\le M_a$, the second by
	$(I-\abs X)^2\le I-\abs X^2$ and $[Q_a,M_a-M_a^2]=0$, and the third by
	$\sum_aP_a=I$.  Minkowski's inequality gives the bound $(3\sqrt\Delta)^2$.
\end{proof}

\begin{lemma}[Consistent measurements]\label{lem:orthonormalization-main-lemma}
	\lean{MIPStarRE.LDT.MakingMeasurementsProjective.SimplifiedOrthogonalization.consistent_measurement_linear_bound, MIPStarRE.LDT.MakingMeasurementsProjective.SimplifiedOrthogonalization.right_consistent_measurement_linear_bound}
	\leanok
	\uses{def:simeq, def:approx_delta}
	Let $\calH_{\mathrm A}$ and $\calH_{\mathrm B}$ be finite-dimensional, let
	$\ket\psi\in\calH_{\mathrm A}\ot\calH_{\mathrm B}$ be a state, and let
	$A$ and $B$ be measurements with the same outcome set such that
	$A_a\ot I\simeq_\zeta I\ot B_a$.  Then there is a projective measurement
	$P$ on $\calH_{\mathrm A}$ with $A_a\ot I\approx_{18\zeta}P_a\ot I$, and
	likewise on $\calH_{\mathrm B}$.
\end{lemma}
\begin{proof}
	\leanok
	\uses{lem:state-dependent-orthogonalization, prop:simeq-for-measurements}
	Let $e=1-\sum_a\bra\psi A_a\ot B_a\ket\psi\le\zeta$.  Direct-sum
	Cauchy--Schwarz and $\sum_aB_a^2\le I$ give
	$(1-e)^2\le\sum_a\bra\psi A_a^2\ot I\ket\psi$.  Apply
	Lemma~\ref{lem:state-dependent-orthogonalization} to the marginal state
	$\varphi(X)=\bra\psi X\ot I\ket\psi$, whose defect is at most
	$2e-e^2\le2\zeta$.
\end{proof}
